\chapter{Integer rectifiability under convergence in mass}\label{sec:integer-current-mass-closure}
Weak compactness of integral currents is a deep theorem, but closure under convergence in mass has a direct proof. We expand \cite{Sim18}, Chapter 6, Lemma 3.17 and its proof, printed pp.~175--176. A countable union of supports need not have locally finite unweighted Hausdorff measure, so the proof uses a sigma-finite carrier and only asserts tangent existence for its locally finite pieces.

\begin{lemma}[A common orientation on a countable union of rectifiable carriers]\label{lem:common-countable-current-carrier}
Let $M_j\subset U\subset\mathbb R^N$ be measurable, countably $m$-rectifiable sets with locally finite $\mathcal H^m$ measure, where $U$ is open and $0\leq m\leq N$. There is a measurable unit simple $m$-vector field $\eta$ on $M=\bigcup_jM_j$, outside a null set, such that $\eta(x)$ orients $T_xM_j$ at almost every point of each $M_j$. The measure $\mathcal H^m\llcorner M$ is sigma-finite, but need not be locally finite. No approximate tangent plane to the entire union is claimed.
\end{lemma}
\begin{proof}
For $m=0$, each carrier is locally finite and their union is countable; take the degree-zero orientation $\eta=1$. All assertions then follow from counting measure. Assume $m\geq1$ for the rest of the proof. Use Borel representatives for the sets, and remove the null exceptional sets where their tangent planes fail to exist. On every pairwise intersection, Corollary \ref{cor:rectifiable-common-tangents} makes the tangent planes agree almost everywhere; their union is locally finite because the two sets are. Remove the countable union of these exceptional sets. Choose a measurable orientation on each $M_j$ by partitioning among its regular $C^1$ charts. Assign $x\in M$ to the first set $M_j$ containing it, and use that set's orientation for $\eta(x)$. The assignment is measurable. At every retained point lying in another $M_k$, tangent agreement makes the assigned vector an orientation of $T_xM_k$ too, possibly opposite to its previously chosen orientation. A countable exhaustion of $U$ by relatively compact open sets, intersected with each $M_j$, gives sigma-finiteness. This argument does not apply a locally finite tangent theorem directly to $M$.
\end{proof}

\begin{lemma}[Mass distance equals the L1 distance of signed multiplicities]\label{lem:current-common-carrier-mass-distance}
Let $T_j=\llbracket M_j,\theta_j,\tau_j\rrbracket$ be locally integer rectifiable currents and use $M,\eta$ from Lemma \ref{lem:common-countable-current-carrier}. Define the signed integer function $a_j$ on $M$ by $a_j\eta=\theta_j\tau_j$ on $M_j$ and $a_j=0$ elsewhere. Then, for each relatively compact open $W\subset U$,
\[
 \mathbf M_W(T_j-T_k)
 =\int_{M\cap W}|a_j-a_k|\,d\mathcal H^m,
 \qquad
 \mathbf M_W(T_j)=\int_{M\cap W}|a_j|\,d\mathcal H^m.
\]
The equalities concern mass on open sets; no boundary mass hypothesis is required.
\end{lemma}
\begin{proof}
The common orientation makes $a_j$ well defined almost everywhere and integer valued, with $|a_j|=\theta_j$ on its carrier. Both currents act by integration against the vector density $a_j\eta$ relative to the same sigma-finite Hausdorff measure. Their difference has density $(a_j-a_k)\eta$. This density is integrable locally since $|a_j-a_k|\leq|a_j|+|a_k|$. Its vector-measure variation is $|a_j-a_k|\mathcal H^m$, because the simple vector $\eta$ has mass norm one. Theorem \ref{thm:current-mass-representation} and Proposition \ref{prop:current-open-mass-variation} identify this variation on $W$ with current mass. The same argument with one coefficient zero proves the second equality. The reference measure itself need not be locally finite: the integrable densities define locally finite vector measures, to which the representation and variation statements apply.
\end{proof}

\begin{theorem}[Completeness of integer rectifiable currents in local mass]\label{thm:integer-rectifiable-local-mass-completeness}
Let $T_j$ be locally integer rectifiable $m$-currents in an open $U\subset\mathbb R^N$, with $0\leq m\leq N$. Suppose they are Cauchy in every local mass seminorm:
\[
 \lim_{j,k\to\infty}\mathbf M_W(T_j-T_k)=0
 \qquad\text{for every open }W\Subset U.
\]
Then there is a locally integer rectifiable current $T$ with $\mathbf M_W(T_j-T)\to0$ for every such $W$. This is \cite{Sim18}, Lemma 3.17. It asserts neither convergence of boundaries in mass nor the closure conclusion under merely weak convergence.
\end{theorem}
\begin{proof}
Use the common carrier and signed coefficients of Lemma \ref{lem:current-common-carrier-mass-distance}. On every relatively compact open set they form a Cauchy sequence in $L^1(\mathcal H^m\llcorner M)$. Completeness of $L^1$ gives limits which agree on overlaps by uniqueness. An increasing countable exhaustion therefore supplies one measurable locally integrable coefficient $a$ on $M$. On each exhaustion set choose an almost-everywhere convergent subsequence of the $L^1$ convergent coefficients. A finite limit of integers is an integer, so $a$ is integer valued almost everywhere. This conclusion is compatible across the exhaustion.

Set $M_+=\{x\in M:|a(x)|>0\}$. Since $|a|\geq1$ there,
\[
 \mathcal H^m(M_+\cap W)\leq\int_{M\cap W}|a|\,d\mathcal H^m<\infty.
\]
It is countably rectifiable as a measurable subset of the countable carrier, and now it has locally finite Hausdorff measure. In dimension zero the local measure bound makes $M_+$ locally finite, and its scalar orientation is immediate. For $m\geq1$, on $M_+\cap M_j$ the common-tangent theorem identifies its tangent with that of $M_j$ outside a null set. Countably many $j$ cover $M_+$, so $\eta$ is tangent to $M_+$ almost everywhere. Consequently $|a|$ and $\operatorname{sign}(a)\eta$ define a locally integer rectifiable current $T$. The mass-distance proof applies to $a_j-a$ and gives $\mathbf M_W(T_j-T)=\int_{M\cap W}|a_j-a|\to0$. The resulting current is unique because mass convergence implies weak convergence on every test. This proof does not assume locally finite Hausdorff measure for the initial countable union and does not invoke Federer--Fleming compactness.
\end{proof}

\begin{corollary}[Completeness when both a current and its boundary converge in mass]\label{cor:integral-current-local-normal-completeness}
Let $0\leq m\leq N$ and let $T_j$ be locally integral $m$-currents in $U$. If both $T_j$ and $\partial T_j$ are Cauchy in every local mass seminorm, there is a locally integral $T$ such that
\[
 \mathbf M_W(T_j-T)+\mathbf M_W(\partial T_j-\partial T)\longrightarrow0
 \quad\text{for every }W\Subset U.
\]
No common compact support is assumed or concluded.
\end{corollary}
\begin{proof}
For $m=0$ apply Theorem \ref{thm:integer-rectifiable-local-mass-completeness} to the currents; all boundaries are zero. For $m\geq1$, the same theorem gives locally integer rectifiable limits $T$ for the currents and $S$ for their boundaries. For a smooth compactly supported $(m-1)$-form $\omega$, choose $W\Subset U$ containing its support. The local mass convergence gives $T_j(d\omega)\to T(d\omega)$ and $\partial T_j(\omega)\to S(\omega)$. Since these two evaluations are equal for every $j$, it follows that $S=\partial T$. Thus $T$ and its boundary are both locally integer rectifiable, and the two already established mass convergences give the conclusion. The same argument without the Cauchy assumption on the boundaries would give no boundary mass convergence; that additional assumption has been used essentially.
\end{proof}
